\section{总结与延伸}

CS329A 的九讲形成了一条完整递进链：从 LLM 与 agent 基础出发，在测试时用采样、搜索、工具和 verifier 提升解题；再把成功轨迹转成 STaR、GRPO、DAPO 等训练信号；随后面对长时程规划、deep research、真实工作评估与开放环境。

最终讲提出四个决定下一阶段进展的研究问题：

\begin{enumerate}
    \item \textbf{如何持续产生新 reasoning？} Multi-agent specialization 用结构化差异维持 synthetic data diversity；
    \item \textbf{如何相信反馈？} DeepSeekMath-V2 用 meta-verification 审计 verifier 的问题分析；
    \item \textbf{如何突破题库？} Absolute Zero 让 proposer 与 solver 共同生成动态 curriculum；
    \item \textbf{如何负担更多计算？} Intelligence per watt 把能力、硬件、能耗和 local/cloud serving 放到同一优化框架。
\end{enumerate}

最值得保留的工程判断是：自我改进不是模型在真空中“自己变聪明”，而是一个由环境、数据、搜索、验证、训练、系统和能源共同约束的闭环。闭环的能力上限由最弱 verifier、最窄数据分布和最昂贵资源共同决定。

因此，一个高质量自我改进 agent 应具备：

\begin{itemize}
    \item 多样 proposal 与可度量的探索覆盖；
    \item 分层、独立、能被反驳的 verifier；
    \item 随能力变化的动态任务课程；
    \item 对慢反馈和 subjective reward 的不确定性建模；
    \item memory 与 parameter learning 的明确边界；
    \item workload-aware 的 local/cloud 路由与能效监控；
    \item 在真实环境和人类目标上的持续外部校准。
\end{itemize}

课程结束，但这一领域变化极快。论文会过时，系统原理仍会复用：扩大搜索空间、改善反馈质量、选择恰当训练信号，并把每单位计算转化为更多可靠智能。

\end{document}
